% GENERATED FILE. Edit canonical lessons, then run npm run generate:books.
% canonical-source-hash: fnv1a64:9578a431c28ec69d
% canonical-lessons: SA-C497-kanyadanam, SA-C497-panigrahanam, SA-C497-vedi, SA-C497-juhoti, SA-C497-saubhagyam

\chapter{Wedding Rites}
\label{ch:sa-wedding-rites}
\hlchaptermodality{497}

\begin{chapteropening}
\textbf{By the end of this chapter:} \emph{I can name the rites of a wedding.}

Complete the last lesson of chapter 497: I can name the rites of a wedding.
\end{chapteropening}

\section[\texorpdfstring{\sk{कन्यादानम्} (kanyādānam)}{कन्यादानम् (kanyādānam)}]{\sk{कन्यादानम्} (kanyādānam) — giving away the bride}

\label{lesson:SA-C497-kanyadanam}



\begin{quote}
Before the new one: say the Sanskrit for a pavilion, then the Sanskrit for the seven steps.
\end{quote}

\subsection*{You'll want to know: \sk{कन्यादानम्}}
\textbf{\sk{कन्यादानम्}} — \emph{kanyādānam} — \textquotedblleft{}giving away the bride\textquotedblright{}.

\begin{quote}
\sk{पिता} \sk{कन्यादानं} \sk{करोति।}
\end{quote}

\begin{quote}
\emph{pitā kanyādānaṃ karoti.}
\end{quote}

\begin{quote}
The father gives away the bride.
\end{quote}

A neuter noun.

\begin{cousinweb}
\textbf{\sk{कन्यादानम्}} (\emph{kanyādānam}) is built from \emph{kanyā}, \textquotedblleft{}a girl, a daughter\textquotedblright{}, and \emph{dāna}, \textquotedblleft{}a gift\textquotedblright{}, from the root √dā of \textbf{\sk{ददाति}}; MW gives \textquotedblleft{}giving a girl in marriage\textquotedblright{} (MW).
\end{cousinweb}

\begin{grammarlens}[title={What you've built}]
One more word for describing what you lived and felt.
\end{grammarlens}

\subsection*{Your turn}
\begin{itemize}
\raggedright
  \item \emph{Say it:} \emph{kanyādānam}
  \item \emph{Say it:} \emph{kanyādānam}, once more
  \item \emph{Say it:} \emph{pitā kanyādānaṃ karoti.} — The father gives away the bride.
  \item \emph{Recall:} say the Sanskrit for a pavilion, then the Sanskrit for the seven steps, then say \emph{kanyādānam} again
\end{itemize}

\begin{tcolorbox}[breakable,colback=teal!4,colframe=teal!35!black,title={Before you move on}]
What does \sk{कन्यादानम्} mean? (\textquotedblleft{}Giving away the bride\textquotedblright{}.) Say it once more.
\end{tcolorbox}
\section[\texorpdfstring{\sk{पाणिग्रहणम्} (pāṇigrahaṇam)}{पाणिग्रहणम् (pāṇigrahaṇam)}]{\sk{पाणिग्रहणम्} (pāṇigrahaṇam) — taking the bride's hand}

\label{lesson:SA-C497-panigrahanam}



\begin{quote}
Before the new one: say the Sanskrit for the seven steps, then the Sanskrit for giving away the bride.
\end{quote}

\subsection*{You'll want to know: \sk{पाणिग्रहणम्}}
\textbf{\sk{पाणिग्रहणम्}} — \emph{pāṇigrahaṇam} — \textquotedblleft{}taking the bride's hand\textquotedblright{}.

\begin{quote}
\sk{पाणिग्रहणम्} \sk{अग्नेः} \sk{समीपे} \sk{भवति।}
\end{quote}

\begin{quote}
\emph{pāṇigrahaṇam agneḥ samīpe bhavati.}
\end{quote}

\begin{quote}
The taking of hands happens beside the fire.
\end{quote}

A neuter noun.

\begin{cousinweb}
\textbf{\sk{पाणिग्रहणम्}} (\emph{pāṇigrahaṇam}) is built from \emph{pāṇi}, \textquotedblleft{}the hand\textquotedblright{}, and \emph{grahaṇa}, \textquotedblleft{}taking\textquotedblright{}, from the root √grah of \textbf{\sk{गृह्णाति}}; MW gives it as \emph{pāṇi-graha}, \textquotedblleft{}taking the bride by the hand, marriage\textquotedblright{} (MW).
\end{cousinweb}

\begin{grammarlens}[title={What you've built}]
One more word for describing what you lived and felt.
\end{grammarlens}

\subsection*{Your turn}
\begin{itemize}
\raggedright
  \item \emph{Say it:} \emph{pāṇigrahaṇam}
  \item \emph{Say it:} \emph{pāṇigrahaṇam}, once more
  \item \emph{Say it:} \emph{pāṇigrahaṇam agneḥ samīpe bhavati.} — The taking of hands happens beside the fire.
  \item \emph{Recall:} say the Sanskrit for the seven steps, then the Sanskrit for giving away the bride, then say \emph{pāṇigrahaṇam} again
\end{itemize}

\begin{tcolorbox}[breakable,colback=teal!4,colframe=teal!35!black,title={Before you move on}]
What does \sk{पाणिग्रहणम्} mean? (\textquotedblleft{}Taking the bride's hand\textquotedblright{}.) Say it once more.
\end{tcolorbox}
\section[\texorpdfstring{\sk{वेदी} (vedī)}{वेदी (vedī)}]{\sk{वेदी} (vedī) — a fire altar}

\label{lesson:SA-C497-vedi}



\begin{quote}
Before the new one: say the Sanskrit for giving away the bride, then the Sanskrit for taking the bride's hand.
\end{quote}

\subsection*{You'll want to know: \sk{वेदी}}
\textbf{\sk{वेदी}} — \emph{vedī} — \textquotedblleft{}a fire altar\textquotedblright{}.

\begin{quote}
\sk{वेद्याम्} \sk{अग्निः} \sk{ज्वलति।}
\end{quote}

\begin{quote}
\emph{vedyām agniḥ jvalati.}
\end{quote}

\begin{quote}
Fire burns on the altar.
\end{quote}

A feminine noun.

\begin{cousinweb}
\textbf{\sk{वेदी}} (\emph{vedī}) MW records \emph{vedi}, later also \emph{vedī}, \textquotedblleft{}an elevated piece of ground serving for a sacrificial altar\textquotedblright{}; MW gives no derivation, so its origin is uncertain (MW).
\end{cousinweb}

\begin{grammarlens}[title={What you've built}]
One more word for describing what you lived and felt.
\end{grammarlens}

\subsection*{Your turn}
\begin{itemize}
\raggedright
  \item \emph{Say it:} \emph{vedī}
  \item \emph{Say it:} \emph{vedī}, once more
  \item \emph{Say it:} \emph{vedyām agniḥ jvalati.} — Fire burns on the altar.
  \item \emph{Recall:} say the Sanskrit for giving away the bride, then the Sanskrit for taking the bride's hand, then say \emph{vedī} again
\end{itemize}

\begin{tcolorbox}[breakable,colback=teal!4,colframe=teal!35!black,title={Before you move on}]
What does \sk{वेदी} mean? (\textquotedblleft{}A fire altar\textquotedblright{}.) Say it once more.
\end{tcolorbox}
\section[\texorpdfstring{\sk{जुहोति} (juhoti)}{जुहोति (juhoti)}]{\sk{जुहोति} (juhoti) — offers into the fire}

\label{lesson:SA-C497-juhoti}



\begin{quote}
Before the new one: say the Sanskrit for taking the bride's hand, then the Sanskrit for a fire altar.
\end{quote}

\subsection*{You'll want to know: \sk{जुहोति}}
\textbf{\sk{जुहोति}} — \emph{juhoti} — \textquotedblleft{}offers into the fire\textquotedblright{}.

\begin{quote}
\sk{वधूः} \sk{अग्नौ} \sk{लाजान्} \sk{जुहोति।}
\end{quote}

\begin{quote}
\emph{vadhūḥ agnau lājān juhoti.}
\end{quote}

\begin{quote}
The bride offers puffed rice into the fire.
\end{quote}

Past (he …): \sk{सः} \sk{अग्नौ} \sk{लाजान्} \sk{हुतवान्}, \emph{saḥ agnau lājān hutavān}.

\begin{cousinweb}
\textbf{\sk{जुहोति}} (\emph{juhoti}) comes from the root √hu, \textquotedblleft{}to offer into the fire\textquotedblright{}, a third-class verb that doubles the root (\emph{ju-ho-ti}); MW's ppp \emph{huta} is \textquotedblleft{}offered in fire\textquotedblright{} (MW).
\end{cousinweb}

\begin{grammarlens}[title={What you've built}]
One more word for describing what you lived and felt.
\end{grammarlens}

\subsection*{Your turn}
\begin{itemize}
\raggedright
  \item \emph{Say it:} \emph{juhoti}
  \item \emph{Say it:} \emph{juhoti}, once more
  \item \emph{Say it:} \emph{vadhūḥ agnau lājān juhoti.} — The bride offers puffed rice into the fire.
  \item \emph{Recall:} say the Sanskrit for taking the bride's hand, then the Sanskrit for a fire altar, then say \emph{juhoti} again
\end{itemize}

\begin{tcolorbox}[breakable,colback=teal!4,colframe=teal!35!black,title={Before you move on}]
What does \sk{जुहोति} mean? (\textquotedblleft{}Offers into the fire\textquotedblright{}.) Say it once more.
\end{tcolorbox}
\section[\texorpdfstring{\sk{सौभाग्यम्} (saubhāgyam)}{सौभाग्यम् (saubhāgyam)}]{\sk{सौभाग्यम्} (saubhāgyam) — good fortune}

\label{lesson:SA-C497-saubhagyam}



\begin{quote}
Before the new one: say the Sanskrit for a fire altar, then the Sanskrit for offers into the fire.
\end{quote}

\subsection*{You'll want to know: \sk{सौभाग्यम्}}
\textbf{\sk{सौभाग्यम्}} — \emph{saubhāgyam} — \textquotedblleft{}good fortune\textquotedblright{}.

\begin{quote}
\sk{सौभाग्यं} \sk{भवतु} \sk{इति} \sk{वृद्धाः} \sk{वदन्ति।}
\end{quote}

\begin{quote}
\emph{saubhāgyaṃ bhavatu iti vṛddhāḥ vadanti.}
\end{quote}

\begin{quote}
\textquotedblleft{}May good fortune be yours,\textquotedblright{} say the elders.
\end{quote}

A neuter noun.

\begin{cousinweb}
\textbf{\sk{सौभाग्यम्}} (\emph{saubhāgyam}) is built on \emph{su-bhaga}, \textquotedblleft{}lucky\textquotedblright{}, with the lengthening of its first vowel; MW gives \textquotedblleft{}good luck, good fortune, happiness (esp. conjugal felicity)\textquotedblright{} (MW).
\end{cousinweb}

\begin{grammarlens}[title={What you've built}]
One more word for describing what you lived and felt.
\end{grammarlens}

\subsection*{Your turn}
\begin{itemize}
\raggedright
  \item \emph{Say it:} \emph{saubhāgyam}
  \item \emph{Say it:} \emph{saubhāgyam}, once more
  \item \emph{Say it:} \emph{saubhāgyaṃ bhavatu iti vṛddhāḥ vadanti.} — \textquotedblleft{}May good fortune be yours,\textquotedblright{} say the elders.
  \item \emph{Recall:} say the Sanskrit for a fire altar, then the Sanskrit for offers into the fire, then say \emph{saubhāgyam} again
\end{itemize}

\begin{tcolorbox}[breakable,colback=teal!4,colframe=teal!35!black,title={Before you move on}]
What does \sk{सौभाग्यम्} mean? (\textquotedblleft{}Good fortune\textquotedblright{}.) Say it once more.
\end{tcolorbox}
